
\فصل{ریسک‌ها، محدودیت‌ها و زمان‌بندی}

\قسمت{ریسک‌ها}

جدول~\ref{جدول:ریسک} ریسک‌های اصلی پروژه به همراه احتمال، تأثیر و راهکار کاهش آن‌ها را نشان می‌دهد.

\begin{table}[ht]
\centering
\caption{ریسک‌های پروژه‌ی ساکنا}
\label{جدول:ریسک}
\begin{tabular}{|>{\raggedleft\arraybackslash}p{3.5cm}|>{\raggedleft\arraybackslash}p{1.5cm}|>{\raggedleft\arraybackslash}p{1.5cm}|>{\raggedleft\arraybackslash}p{5cm}|}
\hline
\rl{\textbf{ریسک}} & \rl{\textbf{احتمال}} & \rl{\textbf{تأثیر}} & \rl{\textbf{راهکار کاهش}} \\
\hline
\rl{محدودیت زمانی پروژه (ددلاین درس)} & \rl{زیاد} & \rl{زیاد} & \rl{تعریف واضح \lr{MVP}، اولویت‌بندی سخت‌گیرانه‌ی آیتم‌ها، کاهش \lr{scope} در صورت نیاز} \\
\hline
\rl{پیچیدگی سیستم پرداخت و کیف پول} & \rl{متوسط} & \rl{زیاد} & \rl{شبیه‌سازی پرداخت به جای اتصال واقعی به درگاه؛ تست جامع تراکنش‌ها} \\
\hline
\rl{تداخل رزرو هم‌زمان} & \rl{متوسط} & \rl{متوسط} & \rl{استفاده از قفل‌گذاری (\lr{Locking}) در سطح پایگاه داده و تراکنش‌های \lr{ACID}} \\
\hline
\rl{آسیب‌پذیری‌های امنیتی} & \rl{کم} & \rl{زیاد} & \rl{استفاده از \lr{JWT}، \lr{RBAC}، اعتبارسنجی ورودی، \lr{HTTPS}؛ مرور کد} \\
\hline
\rl{تغییر نیازمندی‌ها در طول پروژه} & \rl{زیاد} & \rl{متوسط} & \rl{انعطاف‌پذیری چابک؛ تعریف واضح \lr{MVP}؛ ارتباط مداوم با استاد} \\
\hline
\rl{اختلاف در تفسیر نیازمندی‌ها در تیم} & \rl{متوسط} & \rl{متوسط} & \rl{مستندسازی دقیق در سند چشم‌انداز؛ جلسات پلنینگ منظم} \\
\hline
\end{tabular}
\end{table}

\قسمت{محدودیت‌ها}

\شروع{فقرات}
\فقره \textbf{محدودیت زمانی:} پروژه باید در قالب دو ایتریشن با مجموعاً حدود سه هفته به پایان برسد.
\فقره \textbf{محدودیت نیروی انسانی:} تیم چهار نفره است و هر نفر علاوه بر این پروژه، تعهدات درسی دیگری دارد.
\فقره \textbf{محدودیت هزینه:} پروژه بودجه‌ی نقدی ندارد؛ از ابزارها و زیرساخت‌های رایگان یا دانشجویی استفاده می‌شود.
\فقره \textbf{محدودیت فناوری:} فناوری‌ها باید برای همه‌ی اعضای تیم قابل یادگیری و استفاده باشند.
\فقره \textbf{محدودیت محیط آموزشی:} سامانه باید از نظر آموزشی کاربرد متدولوژی \lr{DAD} را نشان دهد.
\پایان{فقرات}

\قسمت{هزینه و زمان‌بندی}

\زیرقسمت{هزینه}

این پروژه یک پروژه‌ی دانشگاهی است و هزینه‌ی نقدی قابل‌توجهی ندارد.
هزینه‌های اصلی عبارتند از:

\شروع{فقرات}
\فقره \textbf{نیروی انسانی:} حدود ۴ نفر × ۳ هفته = ۱۲ نفر-هفته
\فقره \textbf{زیرساخت:} استفاده از \lr{GitHub}  و \lr{Docker} 
\فقره \textbf{ابزارها:} \lr{VS Code}، \lr{Postman}، \lr{Figma}
\پایان{فقرات}

\زیرقسمت{زمان‌بندی}

جدول~\ref{جدول:زمانبندی} برنامه‌ی زمانی پروژه را نشان می‌دهد.

\begin{table}[ht]
\centering
\caption{برنامه‌ی زمانی ایتریشن‌ها}
\label{جدول:زمانبندی}
\begin{tabular}{|>{\raggedleft\arraybackslash}p{2.5cm}|>{\raggedleft\arraybackslash}p{3.5cm}|>{\raggedleft\arraybackslash}p{6cm}|}
\hline
\rl{\textbf{ایتریشن}} & \rl{\textbf{بازه‌ی زمانی}} & \rl{\textbf{اهداف اصلی}} \\
\hline
\rl{فاز ۱} & \rl{خرداد ۲۳ تا 12 تیر} & \rl{تدوین سند چشم‌انداز، سند معماری و سند نیازمندی‌ها و اولویت آنها، راه‌اندازی محیط توسعه، تعریف معماری، پیاده‌سازی احراز هویت، مدیریت واحدها، درخواست‌های خدماتی؛ \lr{UI} اولیه} \\
\hline
\rl{فاز ۲} & \rl{تیر 13 تا 23 مرداد} & \rl{کامل کردن پروژه طبق \lr{CD Lean}} \\
\hline
\end{tabular}
\end{table}

\زیرقسمت{جلسات فاز ۱}
طبق متدولوژی Lean جلسات باید حداقلی باشند ولی برای ایجاد هماهنگی بیشتر و بیشبرد کار در تیم جلسات را به این شکل تشکیل می‌دهیم:
\شروع{فقرات}
\فقره \textbf{جلسه‌ی پلنینگ:} خرداد ۲۳ --- ساعت ۲۲:۰۰
\فقره \textbf{جلسات \lr{daily coordination meeting}:} روزهای فرد (ویس در گروه --- تاپیک دیلی)
\فقره \textbf{جلسه‌ی رترو و دمو:} 7 تیر
\پایان{فقرات}

\زیرقسمت{جلسات فاز ۲}

\شروع{فقرات}
\فقره \textbf{جلسه‌ی پلنینگ:} تیر 13
\فقره \textbf{جلسات \lr{daily coordination meeting}:} شنبه و چهارشنبه
\فقره \textbf{جلسه‌ی رترو و دمو:} 20 تیر و 15 مرداد
\پایان{فقرات}

